\subsection{YOLO11-seg}

YOLO11-seg merupakan varian YOLO11 yang mendukung tugas \textit{instance segmentation}. Model ini digunakan untuk menghasilkan prediksi lokasi objek, kelas, dan \textit{segmentation mask} dalam satu model sehingga proses klasifikasi lauk dan segmentasi tidak memerlukan dua model yang terpisah \cite{khanam2024,qiu2025}.

\subsubsection*{a. Arsitektur YOLO11-seg}

Secara umum, arsitektur YOLO11 terdiri atas tiga bagian utama, yaitu \textit{backbone}, \textit{neck}, dan \textit{head}. \textit{Backbone} mengekstraksi fitur dari citra masukan, \textit{neck} menggabungkan fitur dari beberapa tingkat resolusi, sedangkan \textit{head} menggunakan fitur tersebut untuk menghasilkan prediksi objek \cite{khanam2024}. Pada YOLO11-seg, bagian keluaran juga mencakup informasi yang diperlukan untuk membentuk segmentasi setiap instans \cite{qiu2025}.

Pada bagian \textit{backbone}, citra diproses melalui lapisan konvolusi dan modul utama seperti C3k2, SPPF (\textit{Spatial Pyramid Pooling--Fast}), dan C2PSA. C3k2 digunakan dalam proses ekstraksi fitur, SPPF memperluas informasi konteks melalui operasi \textit{pooling}, sedangkan C2PSA memberikan mekanisme perhatian terhadap fitur yang diproses \cite{khanam2024,qiu2025}.

Fitur dari \textit{backbone} diteruskan ke bagian \textit{neck}. Bagian ini menggabungkan informasi dari beberapa tingkat resolusi melalui operasi seperti \textit{upsampling} dan \textit{concatenation}. Penggabungan fitur lintas skala memungkinkan informasi spasial dari fitur beresolusi lebih tinggi dipadukan dengan informasi semantik dari fitur yang lebih dalam \cite{khanam2024,qiu2025}.

Hasil dari \textit{neck} kemudian diteruskan ke \textit{head}. Pada YOLO11-seg, bagian deteksi menghasilkan prediksi kelas dan posisi \textit{bounding box}, sedangkan bagian segmentasi menghasilkan informasi \textit{mask} untuk setiap kandidat objek. Salah satu mekanisme yang digunakan adalah pembentukan \textit{prototype mask} yang dikombinasikan dengan koefisien \textit{mask} setiap objek untuk menghasilkan \textit{mask} masing-masing instans \cite{qiu2025}. Penerapan YOLO11-seg pada penelitian lain juga menunjukkan penggunaan C3k2, SPPF, C2PSA, penggabungan fitur lintas skala, serta keluaran segmentasi dalam satu arsitektur \cite{syafira2025}.

\subsubsection*{b. Pelatihan YOLO11-seg}

Pelatihan model dapat dimulai dari bobot pralatih agar fitur visual yang telah dipelajari sebelumnya dapat digunakan sebagai titik awal dan kemudian disesuaikan dengan kelas pada dataset penelitian. Pada penelitian ini, bobot YOLO11-seg disesuaikan untuk mengenali sembilan kelas, yaitu nasi, rendang, ayam goreng, perkedel, sayur nangka, daun singkong, telur dadar, telur kari, dan sambal.

Pada proses pelatihan, citra dan \textit{ground truth} diberikan kepada model untuk menghasilkan prediksi lokasi objek, kelas, dan segmentasi. Prediksi tersebut dibandingkan dengan \textit{ground truth} untuk memperoleh nilai \textit{loss}. Nilai \textit{loss} kemudian digunakan dalam proses \textit{backpropagation} untuk memperbarui bobot model secara berulang. Data validasi digunakan untuk memantau performa model selama pelatihan, sedangkan data pengujian digunakan setelah model akhir diperoleh.

Nilai parameter pelatihan, seperti varian YOLO11-seg, ukuran citra masukan, jumlah \textit{epoch}, \textit{batch size}, konfigurasi augmentasi, dan parameter lain yang digunakan pada implementasi dijelaskan pada Bab III.
